\begin{abstract}
Directional failures in language-conditioned manipulation are commonly attributed to weak language grounding. We test this explanation using matched trials that hold the scene, initial state, sampling seed, checkpoint, and controller fixed while switching only the directional word between LEFT and RIGHT. Across 324 pairs spanning eight completed checkpoint--arena rows and two simulated environments, 253 endpoint pairs (78\%) are ordered in the requested direction, showing that the instruction redirects behavior even when task success differs by direction. Scene geometry can determine whether this redirection becomes a successful placement. Reflecting movable-object positions reverses placement advantage in three checkpoints and flips the favored direction for $\pi_{0.5}$: LEFT/RIGHT completion changes from 4/27 and 25/27 in the original layout to 25/27 and 9/27 after reflection. Moving the reference object across seven lateral positions changes the placement gap gradually, while a symmetric scene reduces the gap across checkpoints that pass the endpoint-redirection control and brings DreamZero's observed gap near zero. Canonical response, however, does not guarantee that the physical relation transfers across linguistic forms. For $\pi_{0.5}$, expressing the same goals from the other object's frame collapses LEFT-minus-RIGHT endpoint separation from 23.3 to 0.2\,cm. Binary directional success is therefore not a grounding test by itself: scene geometry can make language-responsive control look like failure, while familiar wording can make surface-form sensitivity look like semantic understanding.
\end{abstract}

\begin{IEEEkeywords}
Vision-language-action models, robot learning, manipulation, language grounding, evaluation methodology.
\end{IEEEkeywords}
